\uuid{zX8O}
\exo7id{4751}
\titre{exo7 4751}
\auteur{quercia}
\organisation{exo7}
\datecreate{2010-03-16}
\isIndication{false}
\isCorrection{false}
\chapitre{Topologie}
\sousChapitre{Topologie des espaces vectoriels normés}
\cpge{mpsi-mp2i.calcul-matriciel-systemes.e, pcsi.determinants.c, ptsi.determinants.a, mpsi-mp2i.determinants.d, pcsi.calcul-matriciel-systemes-lineaires.d, ptsi.calcul-matriciel-systemes-lineaires.d, bcpst1.matrices, mp-mpi.topologie-espaces-normes.a, pc.espaces-vectoriels-normes.a, psi.espaces-vectoriels-normes.a, mp-mpi.topologie-espaces-normes.b, pc.espaces-vectoriels-normes.b, psi.espaces-vectoriels-normes.b, bcpst2.applications-lineaires-matrices.c}
\module{Analyse}
\niveau{L2}
\difficulte{}

\contenu{
\question{
Soit $A \in \mathcal{M}_p(\R)$ quelconque. Montrer qu'il existe une suite de matrices inversibles
convergeant vers $A$.
}
}
